\begin{table*}[p]
\tiny
\centering
\caption{\textbf{Temporal estimands used to characterize local support--engagement coupling.} Main temporal analyses use adjacent turn pairs and short-range event-history models. Coarser before--after summaries are exploratory boundary checks because scaffolded turns are selected by task difficulty, user motivation and conversation trajectory.}
\label{tab:temporal_estimands}
\setlength{\tabcolsep}{4pt}
\renewcommand{\arraystretch}{1.02}
\begin{tabularx}{\textwidth}{>{\raggedright\arraybackslash}p{0.17\textwidth}>{\raggedright\arraybackslash}p{0.25\textwidth}>{\raggedright\arraybackslash}p{0.22\textwidth}>{\raggedright\arraybackslash}X}
\toprule
Estimand & Definition & Reported in main text & Scope \\
\midrule
Overall adjacent-turn lift & Next-turn constructive probability after scaffolded support versus a non-scaffolded reference response, estimated within each task setting. & Figure~\ref{fig:temporal_coupling}a & Local association between support type and immediately following constructive engagement; not a randomized support effect. \\
Prior-state adjacent contrast & Same next-turn contrast, stratified by prior constructive, active and passive user states. & Figure~\ref{fig:temporal_coupling}b & Tests whether local coupling depends on the user's immediately preceding state. \\
Reverse state-conditioned support & Probability that the next assistant turn is scaffolded after constructive, active and passive user turns. & Figure~\ref{fig:temporal_coupling}c & Describes how scaffolding supply varies with the user's current engagement state. \\
Prior-state support-form regression & Adjacent-turn logistic regression for next-turn constructive engagement, adding prior-state $\times$ M1--M6 interactions. & Figure~\ref{fig:temporal_coupling}d and Supplementary Table~\ref{tab:prior_state_support_form_interactions} & Tests whether local coupling depends jointly on prior user state and support form. \\
Short-range event history & Recent-history windows summarize constructive user turns, active user turns and scaffolded assistant responses over the current user--assistant pair and the preceding one to three pairs; a lag-specific logistic model enters current and lagged user-engagement and scaffolded-response indicators together. & Section~2.5 and Supplementary Table~\ref{tab:event_sequence_lagged_scaffolding} & Tests whether the temporal signal reflects a local interactional environment rather than only a scaffolding-only lag or whole-conversation exposure. \\
Around-first-scaffold contrast & Mean constructive engagement after versus before the first scaffolded assistant turn, among conversations containing scaffolded support. & Supplementary Figure~\ref{fig:supp_around_first_scaffolded} & Exploratory boundary check only; the first scaffolded turn may be selected by difficulty or an already developing exchange. \\
\bottomrule
\end{tabularx}
\end{table*}
